\chapter{Fundamentação Teórica} \label{chap:fundamentacao}

% Objetivo: fundamentar teoricamente previsão financeira com TFT, incerteza
% quantílica e avaliação fora da amostra reprodutível.

\section{Previsão de séries temporais financeiras}
A previsão de retornos em mercados financeiros é historicamente tratada como um problema de fronteira entre teoria informacional e evidência empírica. Em sua segunda revisão da hipótese de mercados eficientes, \cite{FAMA1991} consolida a tese de que preços refletem rapidamente a informação pública disponível, de modo que retornos anormais sustentáveis dependem de informação privada, de exposição a fatores de risco não precificados ou de fricções de mercado específicas. O corolário metodológico relevante para este trabalho é direto: previsibilidade estável da média dos retornos é a exceção, não a regra, em ativos líquidos.

Em contraponto, \cite{LO2004} formula a hipótese de mercados adaptativos, que recolhe os mesmos fatos estilizados sob uma lógica evolucionária: eficiência e ineficiência coexistem ao longo do tempo em função de competição entre estratégias, mudanças de regime e adaptação dos agentes. A consequência prática não é a refutação de Fama, mas o reconhecimento de que sinais preditivos --- quando existem --- tendem a ser contextuais, multivariados e temporalmente localizados. Dois fenômenos amplamente documentados reforçam essa leitura: a não estacionariedade de primeira e segunda ordem dos retornos diários e a aglomeração de volatilidade (\textit{volatility clustering}).

Evidência empírica recente recalibra a expectativa de magnitude do sinal. \cite{GU2020}, ao aplicar aprendizado de máquina a apreçamento de ativos, reportam \(R^2\) fora da amostra da ordem de \(0{,}4\%\) no melhor modelo para retornos mensais de ações individuais, medido contra a previsão nula, e observam que médias históricas tipicamente perdem para essa previsão nula; em frequência diária, a previsibilidade da média é ainda menor. \cite{RAPACH2013}, em panorama da literatura sobre previsão de retornos agregados, chegam a magnitudes compatíveis e discutem como pequenos \(R^2\) fora da amostra podem ainda ser economicamente relevantes. Essa literatura ancora duas expectativas centrais para o presente trabalho: ganhos preditivos na média são modestos, e qualquer protocolo que reporte ganhos sistematicamente altos deve ser tratado com suspeita de \textit{leakage} (ver \hyperref[gl:leakage]{Apêndice B}) ou \textit{data snooping} antes da interpretação econômica. Ao mesmo tempo, o valor agregável de um modelo tende a estar na escala e na forma da distribuição condicional --- volatilidade e quantis ---, e não na sua locação.

A combinação de retornos pouco previsíveis, regimes não estacionários e \(R^2\) fora da amostra pequenos torna a robustez metodológica mais informativa do que um número pontual de erro: comparação pareada contra \textit{baselines} explícitos, alinhamento temporal estrito, avaliação quantílica e correção para múltiplos testes passam a ser condição necessária para que a afirmação ``o modelo superou o \textit{baseline}'' carregue significado. O recorte deste trabalho (modelos \textit{single-asset}, AAPL como piloto, horizontes \(h+1\) e \(h+7\)) reflete essa restrição: não se afirma generalização entre ativos ou períodos que a evidência não sustentaria.

\section{Modelagem com aprendizado profundo para séries temporais}
\cite{LIMZOHREN2021} sistematizam o estado da arte em previsão de séries temporais com aprendizado profundo, organizando famílias de arquiteturas (recorrentes, convolucionais, baseadas em atenção e híbridas) e os papéis que podem desempenhar (previsão pontual, multi-horizonte, probabilística e multivariada). Dois pontos da revisão são particularmente pertinentes. Primeiro, modelos profundos oferecem flexibilidade representacional para capturar não linearidades, interações entre covariáveis heterogêneas e dependências temporais de diferentes escalas, atributos relevantes para previsão financeira com indicadores técnicos, sentimento e fundamentos. Segundo, essa flexibilidade vem acompanhada de risco aumentado de sobreajuste e de sensibilidade a \textit{leakage} temporal sutil.

A consequência metodológica reforça a seção anterior: a adoção de aprendizado profundo em previsão financeira não se justifica por premissa de superioridade universal, mas por aderência ao tipo de dado (multivariado, multi-escala, multi-horizonte) e por um protocolo de avaliação que controle explicitamente as fontes de inflação artificial de desempenho. Não basta declarar partições temporais; é preciso controlar \textit{leakage} de \textit{features} (normalizador do alvo ajustado fora do bloco de treino, indicador cuja janela alcance sessões posteriores à decisão --- a janela retrospectiva que termina na própria decisão é causal ---, sentimento agregado sem corte no fechamento da sessão), \textit{leakage} de seleção (hiperparâmetros escolhidos com dados de teste) e \textit{leakage} de alinhamento (interseção temporal incompleta em comparações pareadas). Em séries financeiras, a dependência serial entre observações vizinhas exige ainda separar os blocos de treino e teste por purga e embargo \cite{LOPEZDEPRADO2018}.

A escolha por uma arquitetura profunda neste trabalho é instrumental: serve para representar interações entre famílias de \textit{features} e horizontes em um mesmo modelo, com mecanismos nativos de seleção de variáveis e atenção temporal. Ela não dispensa, mas torna mais crítica, a disciplina experimental de pré-registro, governança de coorte e alinhamento fora da amostra exato.

\section{Temporal Fusion Transformer (TFT)}
O Temporal Fusion Transformer, proposto por \cite{LIMETAL2021TFT}, é um modelo de previsão multi-horizonte que combina componentes recorrentes e de atenção com mecanismos explícitos de seleção de variáveis e saída quantílica. Sua arquitetura organiza a inferência em estágios funcionais: (i) codificação de covariáveis estáticas; (ii) \textit{Variable Selection Networks} (VSN; ver \hyperref[gl:vsn]{Apêndice B}), que atribuem pesos por passo de tempo às variáveis de entrada; (iii) blocos LSTM no codificador e no decodificador, que carregam estado temporal; (iv) atenção \textit{multi-head} interpretável, que captura dependências de longo prazo; e (v) uma camada de saída que emite, para cada horizonte, múltiplos quantis condicionais do alvo em um único passo de inferência.

O modelo distingue formalmente três tipos de entrada: covariáveis estáticas; entradas \textit{observadas}, que só podem ser medidas a cada passo e são desconhecidas de antemão; e entradas \textit{conhecidas}, que podem ser predeterminadas (por exemplo, o dia da semana). Na formulação de \cite{LIMETAL2021TFT}, o alvo e as entradas observadas entram apenas até o instante da decisão \(t\), enquanto as conhecidas entram até \(t+h\). Essa assimetria de indexação é a base formal do controle de \textit{leakage} no modelo: preço, retornos, indicadores, sentimento e fundamentos são necessariamente observados; apenas o calendário é conhecido.

Três características justificam a escolha do TFT como modelo principal. Primeiro, ele integra covariáveis heterogêneas em um único estimador, alinhado às quatro famílias de \textit{features} adotadas (preço, técnico, sentimento, fundamento). Segundo, é nativamente multi-horizonte: um único ajuste, com decodificador de \(h_{\max}\) passos, fornece todos os horizontes, sem modelos independentes por horizonte. Terceiro, sua saída quantílica produz diretamente o objeto probabilístico exigido pelas hipóteses H1 e H2: quantis condicionais por decisão e horizonte, base para a \textit{pinball loss} e para as métricas de calibração. Cabe registrar que o artigo original utiliza apenas três níveis, \(\{0{,}1;\,0{,}5;\,0{,}9\}\); a grade densa adotada neste trabalho não provém de \cite{LIMETAL2021TFT} (ver Seção~\ref{sec:grade}).

A interpretabilidade nativa do TFT, via pesos da VSN e atenção por passo de tempo, é uma das três fontes de evidência para a análise de contribuição de famílias (H3), triangulada com importância por permutação e ablação. Pesos de seleção indicam quais variáveis o modelo usou sob um protocolo, ativo e período --- e não quais variáveis causam o alvo ---, distinção mantida explicitamente em toda a interpretação.

Limites conhecidos. As saídas quantílicas do TFT são transformações lineares do decodificador, treinadas conjuntamente, mas \textbf{sem restrição de não cruzamento}. Assim como em outros regressores quantílicos treinados por nível, o modelo pode produzir \textit{quantile crossing} (ver \hyperref[gl:crossing]{Apêndice B}) e, no extremo, uma grade colapsada em que todos os níveis coincidem \cite{KOENKER2005,CHERNOZHUKOV2010REARRANGEMENT,BONDELL2010,CANNON2018}. Esses fenômenos motivam dois mecanismos distintos adotados neste trabalho: o rearranjo monotônico, parte integral do estimador entregue, e um \textit{gate} de degeneração que invalida linhas colapsadas sem que o rearranjo as mascare.

\section{Predição probabilística e quantílica}
A regressão quantílica, introduzida por \cite{KOENKERBASSETT1978}, é a base formal para estimação de quantis condicionais sem assumir forma paramétrica para a distribuição do erro. Em vez de minimizar o erro quadrático médio, estima-se cada quantil de interesse pela minimização da \textit{pinball loss} (ver \hyperref[gl:pinball]{Apêndice B}), função assimétrica que penaliza desvios em direções opostas com pesos distintos. Para um nível \(\tau \in (0,1)\), valor observado \(y\) e quantil previsto \(\hat{q}\), a \textit{pinball loss} é definida pela Equação~\ref{eq:pinball}:
\begin{equation}
    L_\tau(y, \hat{q}) = \max\left(\,\tau \cdot (y - \hat{q}),\; (\tau - 1) \cdot (y - \hat{q})\,\right).
    \label{eq:pinball}
\end{equation}
A média da \textit{pinball loss} sobre os níveis da grade e sobre as observações fora da amostra é a métrica primária deste trabalho.

A justificativa para tratar a \textit{pinball} como métrica primária, e não apenas medidas de cobertura como o PICP (ver \hyperref[gl:picp]{Apêndice B}), vem da teoria de regras de pontuação. \cite{GNEITING2007} definem regras estritamente próprias, cuja minimização em esperança ocorre se e somente se a previsão reportada coincide com a distribuição verdadeira do alvo; \cite{GNEITING2011} mostra que a \textit{pinball loss} é estritamente consistente para o quantil de nível \(\tau\). Em consequência, reportar quantis honestos é a estratégia ótima para o previsor, e a mesma função serve simultaneamente como perda de treino e como critério de avaliação. Métricas de cobertura agregada, em contraste, podem ser otimizadas com estratégias degeneradas (intervalos artificialmente largos) e por isso não servem como critério primário, ainda que sejam complementos necessários para o diagnóstico de calibração.

Como perfil distributivo complementar, adota-se o \textit{Continuous Ranked Probability Score} (CRPS). Pela representação quantílica de \cite{GNEITINGRANJAN2011}, o CRPS é a integral da \textit{pinball loss} sobre todos os níveis, conforme a Equação~\ref{eq:crps}:
\begin{equation}
    \mathrm{CRPS}(F, y) = 2\int_0^1 L_\alpha\!\left(y, F^{-1}(\alpha)\right)\, d\alpha,
    \label{eq:crps}
\end{equation}
de modo que a \textit{pinball} média sobre uma grade finita é uma aproximação por quadratura de \(\mathrm{CRPS}/2\); quanto mais densa a grade, melhor a aproximação.

\subsection{Grade densa de quantis}\label{sec:grade}
Todos os modelos deste trabalho emitem a mesma grade de níveis \(\mathcal{T}=\{\tau_1<\dots<\tau_K\}\), com \(K\) entre 7 e 9. Não existe teorema de cardinalidade ótima da grade; a escolha apoia-se em três referências convergentes: (i) a perda quantílica da biblioteca utilizada adota sete níveis por padrão, \(\{0{,}02;\,0{,}1;\,0{,}25;\,0{,}5;\,0{,}75;\,0{,}9;\,0{,}98\}\); (ii) a competição M5 \textit{Uncertainty} exigiu nove quantis como suficientes para descrever a distribuição \cite{MAKRIDAKIS2022}; e (iii) a aproximação por quadratura do CRPS (Equação~\ref{eq:crps}) melhora com a densidade da grade. Uma grade comum a todos os modelos torna a \textit{pinball} pareável nível a nível.

\subsection{Calibração}
A avaliação probabilística distingue duas formas de calibração. A primeira é a calibração marginal por quantil, medida pela frequência empírica de observações abaixo do quantil previsto em cada nível \(\tau\), conforme a Equação~\ref{eq:coverage}:
\begin{equation}
    \mathrm{cobertura}_{\tau} = \frac{1}{N}\sum_{i=1}^{N} \mathbb{1}\!\left[\,y_{i} \leq \hat{q}_{\tau,i}\,\right],
    \label{eq:coverage}
\end{equation}
comparada ao nominal \(\tau\); o conjunto dessas coberturas ao longo da grade compõe o \textit{reliability diagram}. A segunda é a calibração intervalar, medida pelo PICP de um intervalo delimitado por um par de níveis \((\tau_l, \tau_h)\), com cobertura nominal \(\tau_h - \tau_l\) (por exemplo, \(0{,}80\) para o par \((0{,}1;\,0{,}9)\)), conforme a Equação~\ref{eq:picp} \cite{KHOSRAVI2011}:
\begin{equation}
    \mathrm{PICP}_{(\tau_l,\tau_h)} = \frac{1}{N}\sum_{i=1}^{N} \mathbb{1}\!\left[\,\hat{q}_{\tau_l,i} \leq y_{i} \leq \hat{q}_{\tau_h,i}\,\right],
    \label{eq:picp}
\end{equation}
em que \(\mathbb{1}[\cdot]\) é a função indicadora e \(N\) é o número de observações fora da amostra. A nitidez é medida pela largura média do intervalo (MPIW), conforme a Equação~\ref{eq:mpiw}:
\begin{equation}
    \mathrm{MPIW}_{(\tau_l,\tau_h)} = \frac{1}{N}\sum_{i=1}^{N} \left(\hat{q}_{\tau_h,i} - \hat{q}_{\tau_l,i}\right),
    \label{eq:mpiw}
\end{equation}
e é sempre analisada conjuntamente com o PICP: intervalo estreito com cobertura baixa indica modelo sistematicamente errado, e não ``preciso''. O \textit{interval score} de \cite{WINKLER1972}, que combina largura e penalidade por não cobertura em um único número próprio \cite{GNEITING2007}, complementa essa leitura.

A cobertura marginal não detecta violações agrupadas no tempo. Para isso, adota-se o teste de cobertura condicional de \cite{CHRISTOFFERSEN1998}, que decompõe a razão de verossimilhança em um componente de cobertura incondicional --- equivalente ao teste de proporção de falhas de \cite{KUPIEC1995} --- e um componente de independência das violações.

\subsection{\textit{Quantile crossing} e rearranjo monotônico}
Quando os níveis de quantil são estimados sem acoplamento, nada garante \(\hat{q}_{\tau_j} \leq \hat{q}_{\tau_k}\) para \(\tau_j < \tau_k\): a curva quantílica prevista pode ser não monótona, o que é incoerente com qualquer distribuição de probabilidade. A correção canônica é o rearranjo monotônico de \cite{CHERNOZHUKOV2010REARRANGEMENT}, que consiste em ordenar os valores da curva estimada. A Proposição~4 desses autores garante que, para qualquer estimador \(\hat{Q}\) de uma curva quantílica verdadeira \(Q_0\), o rearranjado \(\hat{Q}^{*}\) satisfaz
\begin{equation}
    \|\hat{Q}^{*} - Q_0\|_p \;\leq\; \|\hat{Q} - Q_0\|_p, \qquad \forall\, p \in [1, \infty],
    \label{eq:rearranjo}
\end{equation}
em amostra finita e independentemente de como \(\hat{Q}\) foi obtido, com desigualdade estrita (para \(p \in (1,\infty)\)) quando a curva estimada é estritamente decrescente em um conjunto de medida positiva enquanto a verdadeira é estritamente crescente. Ordenar os \(K\) valores da grade é exatamente esse operador na leitura discreta, com pesos iguais por nível; rearranjar, portanto, nunca afasta a previsão da curva verdadeira e tipicamente a aproxima.

Uma ressalva é obrigatória: o rearranjo corrige \textbf{monotonicidade}, não \textbf{calibração}. A Proposição~4 trata de distância à curva verdadeira, não de cobertura empírica; quantis rearranjados podem continuar mal calibrados, o que mantém necessárias a caracterização de calibração (H1) e a referência conformal. Por isso, este trabalho persiste, para cada nível, o valor bruto e o valor rearranjado, de modo que o efeito do rearranjo seja auditável; a variante que alimenta cada métrica confirmatória é fixada no pré-registro. O rearranjo também não trata a grade colapsada, que passa intacta por ele e é tratada por um \textit{gate} de degeneração próprio (Seção~\ref{sec:metricas}).

\subsection{Predição conformal}
A predição conformal oferece intervalos com garantia de cobertura em amostra finita sob a hipótese de permutabilidade dos dados. A regressão quantílica conformalizada (CQR) de \cite{ROMANO2019CQR} ajusta os quantis de um modelo pelo quantil empírico dos escores de não conformidade calculados em um conjunto de calibração disjunto de tudo o que ajustou ou selecionou o modelo. Em séries temporais a permutabilidade não vale, e a cobertura se degrada sob mudança de distribuição \cite{BARBER2023}; por isso, neste trabalho a CQR é usada como \textit{benchmark} comparativo de calibração, com cobertura reportada como \textbf{empírica}, e não como entrega primária.

\subsection{Aplicação a risco}
Quando os quantis previstos servem de base para métricas de risco, a monotonicidade é pré-requisito de definição: o \textit{Value-at-Risk} é um quantil da distribuição de perdas \cite{JORION2007,MCNEIL2005}. Neste trabalho o VaR é reportado de forma descritiva, sobre a variante rearranjada, com \textit{backtest} de exceções pelos testes de \cite{KUPIEC1995} e \cite{CHRISTOFFERSEN1998}.

\section{Avaliação de modelos fora da amostra e testes estatísticos}
Comparar modelos fora da amostra exige, antes de qualquer teste, alinhamento temporal exato. Modelos distintos podem ter previsões sobre coberturas temporais diferentes (por exigências distintas de aquecimento, por exemplo), e calcular médias de erro sobre amostras não idênticas confunde diferença de modelagem com diferença de cobertura. Este trabalho exige interseção exata por (ativo, partição, horizonte, \textit{target timestamp}, nível de quantil), com exatamente uma observação por ponto alinhado, antes de qualquer estatística pareada.

O teste de \cite{DIEBOLDMARIANO1995} (DM; ver \hyperref[gl:dm]{Apêndice B}) opera sobre a série de perdas pareadas entre dois modelos e testa a hipótese nula de igualdade de perda esperada. Dadas as séries de perdas \(L_t^{(1)}\) e \(L_t^{(2)}\), define-se o diferencial \(d_t = L_t^{(1)} - L_t^{(2)}\) e a estatística da Equação~\ref{eq:dm}:
\begin{equation}
    \mathrm{DM} = \frac{\bar{d}}{\sqrt{\widehat{\mathrm{Var}}(\bar{d})}},
    \label{eq:dm}
\end{equation}
em que \(\bar{d}\) é a média amostral de \(d_t\) e \(\widehat{\mathrm{Var}}(\bar{d})\) é um estimador de variância de longo prazo consistente sob heterocedasticidade e autocorrelação (HAC) \cite{NEWEYWEST1987}. Como previsões a \(h\) passos à frente induzem autocorrelação de ordem até \(h-1\) nas perdas, a defasagem do estimador é \(h-1\). Em amostras pequenas o teste original é distorcido; \cite{HARVEY1997} propõem a correção HLN (ver \hyperref[gl:hln]{Apêndice B}), que ajusta a estatística e usa a distribuição \(t\) de Student com \(N-1\) graus de liberdade. Neste trabalho o teste é unilateral (a hipótese alternativa é que o candidato tem perda menor) e aplicado à \textit{pinball loss}.

Comparar o candidato a vários \textit{baselines} infla a taxa de erro do tipo I quando os p-valores não são corrigidos. O ajuste adotado é o procedimento sequencial de \cite{HOLM1979} (ver \hyperref[gl:holm]{Apêndice B}), que controla a \textit{Family-Wise Error Rate} (FWER) e é adequado a uma família pequena e congelada de comparações. A regra de decisão de H2 não exige superioridade contra todos os comparadores: espera-se superar os \textit{baselines} naive, e empatar com os comparadores fortes é resultado informativo.

Quando se deseja, em vez de um vencedor pontual, identificar o conjunto de modelos estatisticamente indistinguíveis do melhor, recorre-se ao \textit{Model Confidence Set} (MCS; ver \hyperref[gl:mcs]{Apêndice B}) de \cite{HANSEN2011MCS}. O procedimento retorna o conjunto de modelos não rejeitados como superiores ao nível \(\alpha\), com \textit{bootstrap} em blocos para respeitar a dependência serial. Pertencer ao MCS (``não rejeitado como inferior'') é evidência complementar ao DM pareado, e divergências entre as duas leituras são elas mesmas informativas.

Por fim, a seleção de configurações usando o mesmo conjunto fora da amostra usado para inferência --- \textit{data snooping} --- faz o ``melhor'' de uma busca ampla parecer significativo por acaso \cite{WHITE2000,ROMANOWOLF2005}. Este trabalho elimina a busca em vez de corrigi-la: um único candidato e comparadores pré-declarados, sem grau de liberdade escolhido após observar o desempenho confirmatório. A governança de coorte (ver \hyperref[gl:coorte]{Apêndice B}) é parte desse protocolo: toda comparação confirmatória é restrita a execuções da mesma coorte, isto é, mesmo ativo, mesmo conjunto de \textit{features}, mesmo horizonte máximo e mesma geometria de partições.

\section{Reprodutibilidade experimental e rastreabilidade}
\cite{PINEAU2021REPRODUCIBILITY}, em relatório sobre o programa de reprodutibilidade da NeurIPS 2019, documentam que uma fração significativa de trabalhos em aprendizado de máquina permanece difícil de reproduzir mesmo com código aberto, por omissões em dados de entrada, sementes e versões de biblioteca, hiperparâmetros e procedimento de partição e seleção de modelo. A variância entre execuções também precisa ser tratada como parte do resultado: \cite{BOUTHILLIER2021} mostram que a amostragem de dados e a inicialização são fontes relevantes de variância em \textit{benchmarks} e recomendam comparar distribuições de desempenho, e não execuções isoladas.

Esse posicionamento ancora a política adotada: as conclusões devem ser reproduzíveis a partir de artefatos persistidos, sem necessidade de re-treino. Operacionalmente, isso implica três elementos. Primeiro, identidade determinística por execução: o identificador de cada \textit{run} é o SHA-256 da representação canônica de nove campos (ativo, \textit{hash} do conjunto de \textit{features}, número do \textit{trial}, \textit{fold}, semente, versão do modelo, assinatura de configuração, assinatura das partições e versão do \textit{pipeline}), de modo que o mesmo experimento sempre produz o mesmo identificador. Segundo, impressões digitais (\textit{fingerprints}) do conjunto de dados, da configuração e das quatro partições temporais asseguram que comparações pareadas ocorrem sobre entradas idênticas. Terceiro, dados brutos e fatos de execução são imutáveis, e toda tabela analítica é derivada e reconstruível a partir deles.

O pré-registro (ver \hyperref[gl:preregistro]{Apêndice B}) completa a política \cite{NOSEK2018}. Ele fixa, antes da rodada confirmatória, hipóteses, candidato, \textit{baselines}, métrica primária, bandas de calibração, variante quantílica, parâmetros dos testes e a regra de decisão. O documento é congelado e \textit{hasheado}, e o \textit{hash} é ancorado aos artefatos confirmatórios, tornando qualquer alteração posterior detectável; isso elimina graus de liberdade que, na ausência de pré-registro, geram \textit{p-hacking} e seleção adversa de \textit{baselines}.

A combinação de identidade determinística, \textit{fingerprints}, dados imutáveis e pré-registro hasheado permite que as conclusões sejam auditadas e reconstruídas a partir dos artefatos persistidos. Essa propriedade não garante que as conclusões sejam corretas, mas garante que sejam falsificáveis no sentido operacional: qualquer revisor pode reexecutar os procedimentos sobre os mesmos artefatos e chegar aos mesmos números.
